\documentclass[11pt,a4paper]{article}
\usepackage[utf8]{inputenc}
\usepackage[T1]{fontenc}
\usepackage[french]{babel}
\usepackage[margin=1.5cm]{geometry}
\usepackage{titlesec}
\usepackage{enumitem}
\usepackage{hyperref}
\usepackage{xcolor}

\definecolor{headercolor}{RGB}{0, 83, 156}

\titleformat{\section}{\large\bfseries\color{headercolor}}{}{0em}{}[\titlerule]
\titlespacing{\section}{0pt}{12pt}{6pt}

\pagestyle{empty}
\setlength{\parindent}{0pt}

\begin{document}

\begin{center}
    {\Huge\textbf{Pierre Dev}}\\[0.4cm]
    {\large Développeur Python}\\[0.3cm]
    06 78 90 12 34 | pierre.dev@email.com | Lyon, France
\end{center}

\section{Expérience Professionnelle}

\textbf{Développeur Python} \hfill \textit{2020 -- Présent}\\
\textit{Entreprise DEF, Lyon}
\begin{itemize}[noitemsep, topsep=3pt]
    \item Développement d'applications web avec Django et Flask
    \item Analyse de données avec Pandas et NumPy
    \item Mise en place de tests unitaires avec Pytest
\end{itemize}

\textbf{Stagiaire Développement} \hfill \textit{2019 -- 2020}\\
\textit{Entreprise GHI, Paris}
\begin{itemize}[noitemsep, topsep=3pt]
    \item Participation au développement d'un projet de reconnaissance d'images avec TensorFlow
    \item Apprendre les bonnes pratiques de développement avec Git
\end{itemize}

\section{Formation}

\textbf{Master Informatique} \hfill \textit{2018 -- 2020}\\
\textit{Université Lyon 1}

\section{Compétences}

\textbf{Langages:} Python, JavaScript, SQL\\
\textbf{Frameworks:} Django, Flask, TensorFlow\\
\textbf{Bases de données:} MySQL, PostgreSQL, MongoDB\\
\textbf{Systèmes d'exploitation:} Windows, Linux, macOS

\end{document}